\chapter{Controlled Primordial Power-Law Sector}

A constant slow-roll parameter gives a solvable inflationary benchmark.  The scalar mode equation and power spectra are used here as a controlled test of the projector discretization.

\allowdisplaybreaks

\section[Constant-epsilon background]{Constant-\(\epsilon\) background}
Assume
\begin{equation}
0<\epsilon<1,
\qquad
q:=\frac{1}{1-\epsilon},
\qquad
\eta<0,
\end{equation}
and
\begin{equation}
\boxed{
a(\eta)=A(-\eta)^{-q}.
}
\end{equation}
Differentiation gives
\begin{align}
a'&=Aq(-\eta)^{-q-1},\\
\mathcal H=\frac{a'}a&=-\frac q\eta,\\
\mathcal H'&=\frac q{\eta^2}.
\end{align}
Hence
\begin{align}
1-\frac{\mathcal H'}{\mathcal H^2}
&=
1-\frac1q
=\epsilon.
\end{align}
Then
\begin{align}
\mathcal H&:=\frac{a'}a=-\frac q\eta,\\
\epsilon_H&:=1-\frac{\mathcal H'}{\mathcal H^2}=\epsilon.
\end{align}
The comoving Hubble radius obeys
\begin{equation}
(aH)^{-1}=\mathcal H^{-1}=-\frac{\eta}{q}.
\end{equation}
For \(N:=\log(a/a_i)\),
\begin{align}
\frac{(aH)^{-1}}{(a_iH_i)^{-1}}
&=\ee^{-(1-\epsilon)N},\\
\frac{|\Omega_K|}{|\Omega_{K,i}|}
&=\ee^{-2(1-\epsilon)N}.
\end{align}
For constant \(\epsilon=-\dd\log H/\dd N\),
\begin{equation}
H/H_i=\ee^{-\epsilon N},
\qquad
a/a_i=\ee^N.
\end{equation}
Therefore \(aH/(a_iH_i)=\ee^{(1-\epsilon)N}\).  Since
\(|\Omega_K|\propto(aH)^{-2}\), the curvature ratio is the square of the comoving-Hubble-radius ratio.
Hence
\begin{equation}
N\geq\frac{\log(S_K^{-1})}{2(1-\epsilon)}
\Longrightarrow
\frac{|\Omega_K|}{|\Omega_{K,i}|}\leq S_K.
\end{equation}

\section{Canonical scalar mode}
For a canonical single field with constant \(\epsilon\),
\begin{equation}
z:=a\sqrt{2\epsilon}\,M_{\mathrm{Pl}},
\qquad
\nu:=q+\frac12=\frac{3-\epsilon}{2(1-\epsilon)}.
\end{equation}
Then
\begin{equation}
\boxed{
\frac{z''}{z}
=\frac{a''}{a}
=\frac{\nu^2-\frac14}{\eta^2}.
}
\end{equation}
Because \(\epsilon\) is constant, \(z\) differs from \(a\) only by a constant factor.  Also,
\begin{align}
\frac{a''}{a}
&=\frac{q(q+1)}{\eta^2},\\
\nu^2-\frac14
&=\left(q+\frac12\right)^2-\frac14\\
&=q(q+1).
\end{align}
The quadratic mode action is
\begin{equation}
S_v^{(2)}
=\frac12\int\dd\eta\dd^3k
\left[
|v_k'|^2-
\left(k^2-\frac{z''}{z}\right)|v_k|^2
\right].
\end{equation}
Varying \(v_k^*\) and integrating the kinetic term by parts gives
\begin{align}
\delta S_v^{(2)}
&=
-\frac12\int
\delta v_k^*
\left[
v_k''+\left(k^2-\frac{z''}{z}\right)v_k
\right]\dd\eta\dd^3k
+\text{c.c.},
\end{align}
so stationarity yields the mode equation.
Thus
\begin{equation}
\boxed{
v_k''+
\left[k^2-\frac{\nu^2-\frac14}{\eta^2}\right]v_k=0.
}
\end{equation}
The Bunch--Davies fixed-background mode is
\begin{equation}
\boxed{
v_k(\eta)
=\frac{\sqrt{-\pi\eta}}2
\ee^{\ii(2\nu+1)\pi/4}
H_\nu^{(1)}(-k\eta).
}
\end{equation}
Its normalization is
\begin{equation}
v_kv_k^{*\prime}-v_k^*v_k'=\ii.
\end{equation}
The vacuum choice and fixed-background canonical quantization are model inputs in this section;  quantized projector geometry is \text{\normalfont not derived here}.

\section{Exact power-law spectra}
For \(a=A(-\eta)^{-q}\), the superhorizon scalar spectrum is
\begin{equation}
\boxed{
\mathcal P_{\mathcal R}(k)
=
\frac{2^{2\nu-4}\Gamma(\nu)^2}
{\pi^3A^2\epsilon M_{\mathrm{Pl}}^2}
 k^{3-2\nu}.
}
\end{equation}
For \(-k\eta\ll1\),
\begin{equation}
H_\nu^{(1)}(x)
\sim
-\frac{\ii\Gamma(\nu)}{\pi}
\left(\frac2x\right)^\nu.
\end{equation}
Thus
\begin{align}
|v_k|^2
&\sim
\frac{2^{2\nu-2}\Gamma(\nu)^2}{\pi}
(-\eta)^{1-2\nu}k^{-2\nu}.
\end{align}
Since \(\mathcal R_k=v_k/z\), \(z=A(-\eta)^{-q}\sqrt{2\epsilon}M_{\mathrm{Pl}}\), and
\(\nu=q+1/2\), all powers of \(\eta\) cancel.  Finally,
\begin{align}
\mathcal P_{\mathcal R}
&=\frac{k^3}{2\pi^2}|\mathcal R_k|^2,
\end{align}
which gives the displayed coefficient and power of \(k\).
The two-polarization tensor spectrum is
\begin{equation}
\boxed{
\mathcal P_T(k)
=
\frac{2^{2\nu}\Gamma(\nu)^2}
{\pi^3A^2M_{\mathrm{Pl}}^2}
 k^{3-2\nu}.
}
\end{equation}
Therefore
\begin{align}
n_s-1&=3-2\nu=-\frac{2\epsilon}{1-\epsilon},\\
n_t&=-\frac{2\epsilon}{1-\epsilon},\\
\alpha_s&:=\frac{\dd n_s}{\dd\log k}=0,\\
\alpha_t&:=\frac{\dd n_t}{\dd\log k}=0,\\
\boxed{r:=\frac{\mathcal P_T}{\mathcal P_{\mathcal R}}=16\epsilon},\\
\boxed{r=-8(1-\epsilon)n_t}.
\end{align}
The spectral index is the logarithmic derivative of the power law:
\begin{equation}
\frac{\dd\log\mathcal P_{\mathcal R}}{\dd\log k}
=3-2\nu.
\end{equation}
Using
\(\nu=(3-\epsilon)/[2(1-\epsilon)]\) gives
\(3-2\nu=-2\epsilon/(1-\epsilon)\).  The scalar and tensor spectra have the same \(k\)-power, so their runnings vanish.  Dividing their normalization constants gives
\begin{equation}
r
=
\frac{2^{2\nu}}{2^{2\nu-4}/\epsilon}
=16\epsilon,
\end{equation}
and substitution of \(n_t\) gives the exact consistency relation.
The first-order slow-roll relation \(r=-8n_t+O(\epsilon^2)\) follows.

\section{Controlled benchmark}
For
\begin{equation}
\epsilon=0.01,
\qquad
N=60,
\end{equation}
\par\noindent
\begin{minipage}[t]{.49\linewidth}\small
The power-law calculation gives
\begin{align}
q&=1.010101,\\
\nu&=1.510101,\\
n_s&=0.979798,\\
n_t&=-0.020202,\\
r&=0.16,\\
\frac{|\Omega_K|}{|\Omega_{K,i}|}
&=2.545749\times10^{-52},\\
\frac{(aH)^{-1}}{(a_iH_i)^{-1}}
&=1.595540\times10^{-26}.
\end{align}
\end{minipage}\hfill
\begin{minipage}[t]{.49\linewidth}\small
The numerical checks are
\begin{align}
\varepsilon_{\epsilon_H}&=0,\\
\varepsilon_{z''/z}&=0,\\
\varepsilon_{n_s}&=1.179612\times10^{-16},\\
\varepsilon_{n_t}&=2.255141\times10^{-16},\\
\varepsilon_r&=2.775558\times10^{-17},\\
\varepsilon_{\mathrm W}&=3.907985\times10^{-14},\\
p_{v,\mathrm{mid}}&=1.964609.
\end{align}
\end{minipage}\par


\section{Open extensions}
\begin{nextsteps}
\item Derive the conditions that select the inflationary attractor from general initial data.
\item Construct a graceful-exit mechanism and a controlled reheating sector.
\item Determine the vacuum from projector refinement instead of fixing it by hand.
\item Quantize the projector geometry together with the scalar perturbations.
\item Extend the calculation to multifield dynamics and the leading non-Gaussian observables.
\end{nextsteps}
